\documentclass[11pt]{article}
\usepackage[T1]{fontenc}
\usepackage{amsmath,amssymb,geometry}
\geometry{margin=1in}
\begin{document}
\begin{center}
{\Large Two multiplicative bases in a group algebra}\\[4pt]
Disproof of Conjecture 00000002758\\[4pt]
ziangni-sys --- October 8, 2026
\end{center}
\paragraph{The actual group algebra.}
Let $G=C_2=\{1,g\}$, $g^2=1$, and let $A=\mathbb Q[G]$ with its
usual convolution multiplication. Its standard group basis is
$B=\{1,g\}$. Define
\[
 e=\frac{1+g}{2},\qquad B'=\{1,e\}.
\]
These are two genuine vector-space bases: for the second one,
$g=2e-1$, and the coordinates of $a+bg$ in $B'$ are $(a-b,2b)$.
The inverse coordinate map is $(s,t)\mapsto s+te$.

\paragraph{Product closure and nonisomorphism.}
The group basis is closed under multiplication since $g^2=1$.
In the second basis,
\[
 e^2=\frac{1+2g+g^2}{4}=\frac{1+g}{2}=e,
 \qquad 1e=e1=e.
\]
Thus $B'$ is also closed under products, exactly as required by
the problem's explicit definition. The coefficient of $g$ in $e$
is $1/2$, whereas it is zero in $1$, so $e\ne1$.
Every element of $B$ has square $1$. A multiplication-preserving
bijection $B\to B'$ would map the identity to the identity and hence
would force every element of $B'$ to have square $1$. This is
impossible for $e$, whose square equals $e\ne1$. Therefore the two
multiplicative bases are not isomorphic as multiplicative structures.
In particular no algebra automorphism of $A$ maps $B$ onto $B'$:
if $\phi(b)=e$ for $b\in B$, then $e^2=\phi(b^2)=\phi(1)=1$,
a contradiction.

\paragraph{The subgroup hypothesis.}
The group $C_2$ has exactly two subgroups, $\{1\}$ and $C_2$.
Indeed any subgroup containing $g$ is the full group, and otherwise
contains only the identity. The two subgroups have orders one and
two, so they are not isomorphic. Thus there are no distinct
isomorphic subgroup pairs, precisely the stated uniqueness hypothesis.
Nevertheless the two multiplicative bases above are nonisomorphic,
which disproves the conjecture's uniqueness assertion.

\paragraph{Scope of isomorphism.}
The isomorphisms here preserve the basis multiplication, as required
for a meaningful assertion of uniqueness of multiplicative bases.
All vector-space bases of the same size admit linear bijections,
so forgetting multiplication would make the proposed uniqueness
condition vacuous. The source does not require basis elements to
be units; imposing that extra requirement would exclude the
idempotent basis and would be a different claim.

\paragraph{Formal certificate.}
Lean uses the actual monoid algebra of the multiplicative group of
$\mathbb Z/2\mathbb Z$, with its genuine convolution product.
Both bases are constructed from explicit invertible coordinate
maps. Product closure, idempotence and the impossibility of an
algebra automorphism between the bases are certified, as is the
absence of distinct isomorphic subgroups. Final axiom audits use
only standard Lean and Mathlib axioms.
\end{document}
